%% kapa: exam=1 round=23 year=2012 period=2 form=A subject=accounting subject_ko=회계학 seq=6 number=6
%% source: Q-Net 감정평가사 기출문제 (cst00309) artlSeq=5052580 file=raw/1008/5052580/2012년 제23회 감정평가사 제1차시험 기출문제(2교시 A형).pdf
%% conversion: kordoc 4.24.51 -> tools/gen1.py (auto); review=proofread-sonnet; flags=none
\begin{question}{6}
(주)한국은 20×1년 12월 31일에 3개의 사업부 중 사업부 A를 매각하기로 결정하였다. (주)한국과 사업부 A의 20×1년 12월 31일 자산과 부채의 장부금액 및 20×1년의 수익과 비용 발생금액은 다음과 같다. 사업부 A의 자산은 모두 비유동자산으로 순공정가치가 ￦1,800이며, 부채의 순공정가치는 장부금액과 동일하다.
\begin{center}
\small
\begin{tabular}{|c r r|}
\hline
구 분 & \uline{(주)한국} & \uline{사업부 A} \\
자 산 & ￦10,000 & ￦2,000 \\
부 채 & 6,000 & 1,500 \\
자 본 & 4,000 & － \\
수 익 & 8,000 & 1,000 \\
비 용 & 6,000 & 700 \\
\hline
\end{tabular}
\end{center}
(주)한국의 20×1년 포괄손익계산서에 표시될 계속영업손익과 중단영업손익은? \par
\begin{choices}
  \item 계속영업이익 ￦2,000 중단영업이익 ￦100
  \item 계속영업이익 ￦1,700 중단영업이익 ￦100
  \item 계속영업이익 ￦1,800 중단영업이익 ￦100
  \item 계속영업이익 ￦1,700 중단영업이익 ￦300
  \item 계속영업이익 ￦2,000 중단영업이익 ￦300
\end{choices}
\end{question}
\answer{2}
